% =============================================================================
%  Bidirectional Vietnamese Dialect Transfer
%  Capstone report -- IT4772E Natural Language Processing, HUST 2025-2
%
%  Compilation: pdflatex -> bibtex -> pdflatex -> pdflatex
%  (or: latexmk -pdf main.tex). UTF-8 input + lmodern renders Vietnamese
%  diacritics (e.g. "An~h di mo^ ro^`i?") correctly under pdflatex.
%  See README_report.md for details and the expected figure list.
% =============================================================================
\documentclass[11pt]{article}

% --- Encoding / fonts -------------------------------------------------------
% Compile with a Unicode engine: XeLaTeX/LuaLaTeX or `tectonic main.tex`.
% fontspec + the default Latin Modern OpenType font renders all Vietnamese
% precomposed glyphs (e.g. "Ảnh đi mô rồi?", "ồ", "ổ") natively from the UTF-8
% source -- no inputenc/fontenc/babel needed.
\usepackage{fontspec}
\usepackage{float}
\usepackage{microtype}
% The SentencePiece word-boundary marker U+2581 ("▁") is absent from Latin Modern
% Mono; pull it from DejaVu Sans Mono, which has the glyph, only for that symbol.
\newfontfamily\dejavumono{DejaVu Sans Mono}
\newcommand{\spm}{{\dejavumono\symbol{"2581}}}

% --- Math / tables / figures ------------------------------------------------
\usepackage{amsmath}
\usepackage{amssymb}
\usepackage{booktabs}
\usepackage{multirow}
\usepackage{array}
\usepackage{graphicx}
\usepackage{subcaption}
\usepackage{xcolor}
\usepackage{enumitem}
\usepackage[round]{natbib}   % provides \citep / \citet used throughout

% --- Page geometry / spacing ------------------------------------------------
\usepackage[a4paper,margin=1in]{geometry}
\usepackage{titlesec}
\usepackage{caption}
\captionsetup{font=small,labelfont=bf}

% --- Hyperlinks (load last among these) -------------------------------------
\usepackage[colorlinks=true,linkcolor=black,citecolor=blue,urlcolor=blue]{hyperref}
\usepackage{microtype}
\setlength{\emergencystretch}{3em}

\newcommand{\ra}{\(\rightarrow\)\allowbreak{}}

% --- Custom macros ----------------------------------------------------------
% \TODO{...} : a bright-red placeholder for numbers / text to be filled once
%              experiments finish. \result{...} : a neutral slot for a final
%              number (kept distinct so a search-and-replace can find each).
\newcommand{\TODO}[1]{\textcolor{red}{\textbf{[TODO: #1]}}}
\newcommand{\result}[1]{\TODO{#1}}      % alias: every result is currently a TODO
\newcommand{\figures}{figures} % figures generated by a separate step

% prefix token rendering helper
\newcommand{\tprefix}[1]{\texttt{#1}}

\title{\textbf{Bidirectional Vietnamese Dialect Transfer:\\
Region-Conditioned Standard-to-Dialect Generation with\\
Multi-Task Seq2Seq and Sampling Back-Translation}}

\author{%
  Capstone Project --- IT4772E Natural Language Processing\\
  Hanoi University of Science and Technology (HUST), Semester 2025-2
}
\date{}

\begin{document}

\begin{titlepage}
  \newgeometry{top=2.2cm,bottom=2.2cm,left=2.5cm,right=2.5cm}
  \thispagestyle{empty}

  \begin{center}
    {\Large \textbf{HANOI UNIVERSITY OF SCIENCE AND TECHNOLOGY}}\\[0.25cm]
    {\large School of Information and Communication Technology}\\[1.2cm]

    \rule{\linewidth}{0.8pt}\\[0.35cm]
    {\Huge \textbf{CAPSTONE PROJECT}}\\[0.25cm]
    \rule{\linewidth}{0.8pt}\\[1.2cm]

    {\Large \textbf{IT4772E Natural Language Processing}}\\[0.35cm]
    {\large Semester 2025--2}\\[1.4cm]

    {\LARGE \textbf{Bidirectional Vietnamese Dialect Transfer}}\\[0.35cm]
    {\Large Region-Conditioned Standard-to-Dialect Generation}\\[0.2cm]

    \vspace{36pt}
    \begin{minipage}{0.78\linewidth}
      \centering
      {\large \textbf{Team Members}}\\[0.35cm]

      \begin{tabular}{@{}l c@{}}
        \toprule
        \textbf{Name} & \textbf{Student ID} \\
        \midrule
        Chu Xuân Minh       & 20235527 \\
        Trần Vũ Mạnh Đức    & 20235487 \\
        Nguyễn Tiến Phát    & 20235544 \\
        Nguyễn Bình Anh     & 20235470 \\
        \bottomrule
      \end{tabular}
    \end{minipage}

    \vfill

    {\large Hanoi, 2026}
  \end{center}

  \restoregeometry
\end{titlepage}



\newpage
% =============================================================================

% --- Headline numbers are placeholders until experiments finish. ------------
\begin{abstract}
Vietnamese dialectal variation is prominent in everyday communication, especially
in social-media text, but most existing NLP work treats it primarily as a source
of noise to be normalized into standard Vietnamese. In this work, we study
Vietnamese dialect transfer in both directions. We formulate region-conditioned
standard-to-dialect generation as a complementary task to dialect-to-standard
normalization, targeting the Northern, Central, and Southern dialect regions
within a unified sequence-to-sequence framework. We fine-tune a single
BARTpho-syllable model with task-prefix conditioning to perform three related
mappings: dialect-to-standard normalization, lexical normalization, and
region-conditioned standard-to-dialect generation. To build the generation data,
we reverse the ViDia2Std dialect--standard pairs while retaining the original
human-written dialect sentence as the target. Since the resulting data is
strongly region-imbalanced, we further apply region rebalancing and
sampling-based back-translation for the lower-resource regions. Because
standard-to-dialect generation is inherently one-to-many, we evaluate it with
both conventional overlap metrics and transfer-oriented metrics, including
Dialect Feature Recall, Dialectal Edit-Recall, and copy-rate. The experiments
show that the multi-task model remains competitive on normalization while
learning meaningful dialectal substitutions in the reverse direction. The results
also reveal persistent challenges, particularly for low-resource regions and for
automatic evaluation with single-reference dialectal targets. Overall, this work
provides a practical framework for bidirectional Vietnamese dialect transfer and
establishes region-conditioned dialect generation as a measurable task for future
study.
\end{abstract}

% =============================================================================
\section{Introduction}
\label{sec:intro}

Vietnamese is spoken by roughly one hundred million people across three broad
dialect regions---Northern (centred on Ha Noi), Central (e.g.\ Nghe An, Ha Tinh,
Hue), and Southern (centred on Ho Chi Minh City). These dialects differ in
lexical choice, function words, pronouns, and orthographic realization of
colloquial speech. A Central speaker may write \emph{``Anh đi mô rồi?''}
where a standard speaker writes \emph{``Anh đi đâu rồi?''}
(``Where did he go?''); the Central marker \emph{mô} (``where'') and the Southern
\emph{hông} / \emph{hổng} (``not'') have no place in formal written Vietnamese but
are pervasive on social media.

Almost all prior computational treatment of this variation is
\emph{normalizing}: it maps dialectal or noisy text onto a standard target. The
recent ViDia2Std benchmark \citep{ta2026vidia2std} and the Arabic NADI shared
tasks \citep{nadi2024} both frame the problem as dialect$\rightarrow$standard.
The reverse direction---generating idiomatic regional text \emph{from} a standard
sentence---has, to our knowledge, no published treatment for Vietnamese. Yet this
direction is valuable for data augmentation, for controllable and localized text
generation, for stylistic personas in dialogue systems, and for studying the
linguistic structure of the dialects themselves.

This report describes a capstone system that addresses both directions inside a
\emph{single} model. Our contributions are:

\begin{enumerate}[leftmargin=1.4em,itemsep=2pt,topsep=2pt]
  \item \textbf{A novel task: region-conditioned standard$\rightarrow$dialect
    generation} (Feature~C) for Northern, Central, and Southern Vietnamese,
    realized by reversing the ViDia2Std parallel corpus and conditioning a
    seq2seq model on a region-specific task prefix.
  \item \textbf{A unified multi-task model}: one BARTpho-syllable model performs
    dialect$\rightarrow$standard normalization (A), region-conditioned
    standard$\rightarrow$dialect generation (C), and ViLexNorm lexical
    normalization (B), all selected by a textual task prefix.
  \item \textbf{An evaluation suite appropriate to one-to-many generation}: we
    argue that single-reference BLEU under-rewards valid stylistic diversity
    \citep{edunov2018understanding} and introduce Dialect Feature Recall (DFR), a
    marker-independent Dialectal Edit-Recall, and a marker-based region-accuracy
    measure, complemented by a human-evaluation protocol.
  \item \textbf{Data strategy for severe regional imbalance}: oversampling plus
    sampling-based back-translation \citep{edunov2018understanding} that keeps
    \emph{real} dialect as the target and synthesizes standard sources, directed
    at the data-starved Northern and Southern regions.
\end{enumerate}

On the well-studied normalization direction our syllable-level model is
competitive with much larger systems (Section~\ref{sec:results}), which we
present as a positioning result rather than a claim of state of the art.

% =============================================================================
\section{Related Work}
\label{sec:related}

\paragraph{Vietnamese pretrained seq2seq models.}
Two encoder--decoder backbones dominate Vietnamese generation. BARTpho
\citep{tran2022bartpho} adapts the BART/mBART denoising objective
\citep{lewis2020bart,liu2020mbart} to Vietnamese and is released in two
tokenization variants: \emph{BARTpho-syllable}, which segments text at the
syllable level (the natural orthographic unit of Vietnamese, since white space
separates syllables rather than words), and \emph{BARTpho-word}, which first
applies automatic word segmentation. ViT5 \citep{phan2022vit5} is a Vietnamese
T5 \citep{raffel2020t5} trained with a span-corruption objective. Both are
strong for normalization and summarization; we adopt BARTpho-syllable because
syllable-level tokenization is robust to the out-of-vocabulary, non-standard
spellings that pervade dialectal and social-media text, where a word segmenter
trained on standard Vietnamese is unreliable. For dialect identification we note
PhoBERT \citep{nguyen2020phobert}, a RoBERTa-style Vietnamese encoder, as a
natural classifier backbone, and VnCoreNLP \citep{vu2018vncorenlp} for
segmentation and POS analysis.

\paragraph{Dialect-to-standard normalization.}
ViDia2Std \citep{ta2026vidia2std} is the first large parallel corpus of
Vietnamese dialect/standard sentence pairs, collected from public Facebook
comments across all 63 provinces and region-labelled by native annotators. The
authors frame the task as monolingual normalization (dialect$\rightarrow$standard)
and benchmark a range of seq2seq models, with mBART-large-50 the strongest. The
NADI shared tasks for Arabic \citep{nadi2024} similarly target dialect
identification and dialect$\rightarrow$MSA normalization. ViLexNorm
\citep{nguyen2024vilexnorm} addresses the adjacent problem of \emph{lexical}
normalization---mapping teencode and abbreviated social-media spellings to their
standard forms---and VSEC \citep{nguyen2021vsec} addresses Vietnamese spelling
correction. All of these are \emph{many-to-one} normalization tasks. We
incorporate the normalization direction (Features A and B) but argue that the
inverse, generative direction is both unsolved and qualitatively different.

\paragraph{Text style transfer and dialect generation.}
The reverse direction is a form of text style transfer. Non-parallel methods
learn to disentangle content from style with adversarial cross-alignment
\citep{shen2017style} or via unsupervised-MT machinery
\citep{lample2018unsupervised}. Closest to our work, \citet{le2023vietnamese}
study Central$\leftrightarrow$Northern Vietnamese dialect text transfer, but
restrict themselves to two regions and a pairwise transfer setting. We instead
condition a single model on three regions and treat the standard variety as a
shared pivot, so that generation and normalization share one parameter set and
one tokenizer. To our knowledge no prior work performs region-conditioned
standard$\rightarrow$dialect generation across all three Vietnamese regions.

\paragraph{Back-translation and monolingual data.}
Back-translation \citep{sennrich2016improving} augments a translation system by
pairing real target-side sentences with machine-generated source-side
sentences. \citet{edunov2018understanding} show that \emph{sampling} (or
noised-beam) back-translation outperforms pure beam search, because greedy/beam
decoding produces overly regular synthetic sources that under-represent the true
data distribution and induce mode collapse; they also caution that BLEU
correlates poorly with the diversity and quality that back-translation actually
adds. We follow this prescription directly: we keep \emph{real} dialect
sentences as targets and \emph{sample} synthetic standard sources from a strong
dialect$\rightarrow$standard model, concentrating the augmentation on the
data-starved Northern and Southern regions.

\paragraph{Evaluation of generation.}
Corpus BLEU \citep{papineni2002bleu,post2018sacrebleu}, chrF
\citep{popovic2015chrf}, ROUGE-L \citep{lin2004rouge}, METEOR
\citep{banerjee2005meteor}, and the embedding-based BERTScore
\citep{zhang2020bertscore} are the standard automatic metrics. None of them is
designed to reward producing \emph{a} valid dialectal realization among many; we
therefore complement them with task-specific measures (Section~\ref{sec:metrics}).

% =============================================================================
\section{Background and NLP Fundamentals}
\label{sec:background}

This section develops the modelling and evaluation concepts the system relies
on, in the order they appear in the method.

\subsection{Tokenization for Vietnamese: syllables vs.\ words vs.\ subwords}
\label{sec:bg-tok}

Written Vietnamese delimits \emph{syllables} (\emph{ti\'{e}ng}) with white space;
a ``word'' may comprise one or several syllables (e.g.\ \emph{sinh vi\^{e}n},
``student'', is two syllables but one word). A model can therefore tokenize at
three granularities: (i) automatic \emph{word} segmentation, as performed by
VnCoreNLP \citep{vu2018vncorenlp} and used by BARTpho-word and PhoBERT
\citep{nguyen2020phobert}; (ii) raw \emph{syllables}, used by BARTpho-syllable
\citep{tran2022bartpho}; and (iii) statistical \emph{subword} units such as BPE
\citep{sennrich2016bpe} applied on top of either. Word segmentation injects
useful structure on clean text, but it is brittle on dialectal and social-media
input, where non-standard spellings (\emph{kh\^{o}ng}\,$\to$\,\emph{h\^{o}ng}),
elongations, and region-specific function words are out-of-vocabulary for a
segmenter trained on standard Vietnamese. A segmentation error then propagates to
every downstream token. Syllable-level tokenization sidesteps this: every
white-space-delimited unit is a valid token, the vocabulary is small and stable,
and rare dialectal forms degrade gracefully into shorter, still-meaningful units.
For a task whose entire input distribution is dialectal/noisy text, this
robustness is decisive, which is why we adopt BARTpho-syllable.

\subsection{Conditional sequence-to-sequence generation}
\label{sec:bg-seq2seq}

All three tasks are cast as conditional generation. Given a source token sequence
$x = (x_1,\dots,x_m)$ and a discrete \emph{condition} $c$ (here a task/region
prefix), an encoder--decoder Transformer \citep{vaswani2017attention} models the
target $y = (y_1,\dots,y_n)$ autoregressively:
\begin{equation}
  P_\theta(y \mid x, c) \;=\; \prod_{t=1}^{n} P_\theta\!\left(y_t \mid y_{<t},\, x,\, c\right).
  \label{eq:seq2seq}
\end{equation}
We realize the condition $c$ \emph{textually}, by prepending a task prefix to the
source so that the input is the string ``\texttt{\{task\}: \{source\}}''. This
keeps the architecture unchanged---no new embeddings or heads---and lets a single
set of parameters serve every task, with the prefix steering generation.

\subsection{Training objective: teacher forcing, cross-entropy, label smoothing}
\label{sec:bg-train}

The model is trained by \emph{teacher forcing}: at decoder step $t$ the
ground-truth prefix $y_{<t}$ is fed in regardless of the model's own predictions,
which makes the per-step targets independent and the loss fully parallelizable
over $t$. The objective is token-level cross-entropy (negative log-likelihood),
\begin{equation}
  \mathcal{L}(\theta) \;=\; -\sum_{t=1}^{n} \log P_\theta\!\left(y_t \mid y_{<t}, x, c\right).
  \label{eq:ce}
\end{equation}
Minimizing Eq.~\ref{eq:ce} drives the model toward a one-hot target distribution,
which can over-confidently penalize the many \emph{other} valid continuations
that exist in a one-to-many task. \emph{Label smoothing}
\citep{szegedy2016labelsmoothing} mixes the one-hot target with a uniform
distribution of mass $\epsilon$, replacing the target probability of the gold
token with $1-\epsilon$ and distributing $\epsilon$ over the vocabulary. This
regularizes the model, improves calibration, and is especially appropriate when
several outputs are acceptable---as in standard$\rightarrow$dialect generation.

\subsection{Decoding: beam search vs.\ sampling}
\label{sec:bg-decode}

At inference we must search Eq.~\ref{eq:seq2seq} for an output. \emph{Beam
search} keeps the $k$ highest-probability partial hypotheses at each step and
returns the most probable complete sequence; it is the right choice when the task
is essentially deterministic (one correct answer), as in normalization.
\emph{Sampling} draws $y_t \sim P_\theta(\cdot \mid y_{<t}, x, c)$, optionally
restricted to the top-$k$ tokens or to the smallest set whose cumulative
probability exceeds $p$ (nucleus / top-$p$ sampling) and rescaled by a temperature
$T$. Sampling trades a little adequacy for diversity. The distinction matters
twice in this work: (i) for \emph{generating} dialect we want some diversity, and
(ii) for \emph{back-translation} we deliberately sample so that synthetic sources
are not the artificially regular outputs of beam search
\citep{edunov2018understanding}.

\subsection{Back-translation}
\label{sec:bg-bt}

Back-translation \citep{sennrich2016improving,edunov2018understanding} exploits
abundant target-side data when parallel data is scarce. To build augmented pairs
for a source$\to$target direction, one runs a \emph{reverse} model
(target$\to$source) over real target sentences to produce synthetic sources, then
trains on the (synthetic-source, real-target) pairs. Because the target is real,
the model never learns to produce machine artefacts; because the source is only an
input, its imperfections are tolerable. For our standard$\rightarrow$dialect task
the ``target'' is real dialect, so we run the dialect$\rightarrow$standard model
to hallucinate standard sources---and, following
\citet{edunov2018understanding}, we \emph{sample} rather than beam-decode them.

\subsection{The one-to-many problem and its evaluation challenge}
\label{sec:bg-one2many}

Normalization is many-to-one: many dialectal spellings collapse to one standard
form, so a single reference is (nearly) the only correct answer and reference-based
overlap metrics are meaningful. Generation is the inverse, one-to-many: one
standard sentence has many equally valid dialectal realizations (different markers,
pronouns, and orthographic choices), and the corpus supplies only one. A model that
produces a perfectly idiomatic alternative to the reference is penalized by BLEU
exactly as a wrong answer would be, and---because dialectal rewrites change only a
few tokens---a model that simply \emph{copies} the standard input can score
deceptively high on overlap metrics while performing no transfer at all. This is
the concrete form of the warning in \citet{edunov2018understanding} that
n-gram overlap correlates poorly with stylistic quality and diversity. It
motivates the dedicated metrics of Section~\ref{sec:metrics}: measures of whether
the \emph{right kind} of change was made, independent of matching one particular
reference, plus an explicit copy-rate to expose under-generation.

% =============================================================================
\section{Datasets}
\label{sec:data}

\paragraph{ViDia2Std (Features A and C).}
ViDia2Std \citep{ta2026vidia2std} provides 13{,}657 dialect/standard sentence
pairs mined from public Facebook comments across all 63 Vietnamese provinces and
labelled by native annotators with one of three regions (Northern, Central,
Southern). We use the corpus in \emph{both} directions: read forward
(dialect$\rightarrow$standard) it supplies Feature~A; read in reverse
(standard$\rightarrow$dialect) the same pairs, retagged with a region-specific
prefix, supply the novel Feature~C. The corpus is severely region-imbalanced
(Central $\approx$66\%, Southern $\approx$27\%, Northern $\approx$7\%;
Figure~\ref{fig:imbalance}), which directly motivates the rebalancing and
back-translation of Section~\ref{sec:method}.

\paragraph{ViLexNorm (Feature B).}
ViLexNorm \citep{nguyen2024vilexnorm} provides 10{,}467 lexical-normalization
pairs (teencode/social-media text $\rightarrow$ standard Vietnamese) with an
8:1:1 train/dev/test split. Including it as a third task exposes the shared
encoder to a broad range of non-standard spellings, which we expect to help the
dialect tasks that face the same orthographic noise.

% \paragraph{Unified format.}
% Every example is stored as a JSON record with fields \texttt{task},
% \texttt{region}, \texttt{source}, and \texttt{target}, and is fed to the model as
% the prefixed string ``\texttt{\{task\}: \{source\}}''. The four canonical task
% labels are Dialect-to-standard, Northern standard-to-dialect,
% Central standard-to-dialect, Southern standard-to-dialect, plus
% Lexical normalization. Table~\ref{tab:datastats} summarizes the per-task split sizes;
% Figure~\ref{fig:composition} visualizes the overall composition.

% \begin{table}[t]
%   \centering
%   \small
%   \caption{Per-task example counts in the processed splits. ``\,---\,'' marks
%   regions that do not apply to a task. Counts are read off the processed data.}
%   \label{tab:datastats}
%   \begin{tabular}{l r r r}
%     \toprule
%     \textbf{Task (prefix)} & \textbf{Train} & \textbf{Dev} & \textbf{Test} \\
%     \midrule
%     Dialect-to-standard (Feature A)          & 10{,}844 & 1{,}181 & 1{,}597 \\
%     Central standard-to-dialect (Feature C) &  7{,}196 &   767 & 1{,}066 \\
%     Southern standard-to-dialect (Feature C)&  2{,}893 &   319 &   330 \\
%     Northern standard-to-dialect (Feature C)&    755 &    95 &   201 \\
%     Lexical normalization (Feature B)              &  8{,}372 & 1{,}050 & 1{,}045 \\
%     \midrule
%     \textbf{Total}                             & 30{,}060 & 3{,}412 & 4{,}239 \\
%     \bottomrule
%   \end{tabular}
% \end{table}


\paragraph{Unified format.}
All datasets are converted to a common text-to-text format: each example has a
task condition, a source sentence, and a target sentence. The condition selects
one of the five mappings: dialect-to-standard normalization, lexical
normalization, or standard-to-dialect generation for one of the three regions.
In implementation, this condition is realized as a short prefix attached to the
source sentence. Table~\ref{tab:datastats} reports the resulting split sizes, and
Figure~\ref{fig:composition} shows the overall task composition.

\begin{table}[H]
  \centering
  \small
  \caption{Per-task example counts in the processed splits.}
  \label{tab:datastats}
  \begin{tabular}{@{}l r r r@{}}
    \toprule
    \textbf{Task} & \textbf{Train} & \textbf{Dev} & \textbf{Test} \\
    \midrule
    Dialect-to-standard normalization       & 10{,}844 & 1{,}181 & 1{,}597 \\
    Central standard-to-dialect generation  &  7{,}196 &   767 & 1{,}066 \\
    Southern standard-to-dialect generation &  2{,}893 &   319 &   330 \\
    Northern standard-to-dialect generation &    755 &    95 &   201 \\
    Lexical normalization                   &  8{,}372 & 1{,}050 & 1{,}045 \\
    \midrule
    \textbf{Total}                          & 30{,}060 & 3{,}412 & 4{,}239 \\
    \bottomrule
  \end{tabular}
\end{table}



\begin{figure}[t]
  \centering
  \includegraphics[width=0.7\linewidth]{\figures/region_imbalance.pdf}
  \caption{Regional imbalance in ViDia2Std: Central dialect dominates while
  Northern is severely under-represented, motivating oversampling and targeted
  back-translation.  }
  \label{fig:imbalance}
\end{figure}

\begin{figure}[t!]
  \centering
  \includegraphics[width=0.8\linewidth]{\figures/dataset_composition.pdf}
  \caption{Composition of the unified multi-task training set across the five
  task prefixes (three regions of Feature~C, plus Features~A and~B).
  }
  \label{fig:composition}
\end{figure}

% =============================================================================
% \section{Method}
% \label{sec:method}

% \subsection{Architecture and multi-task task-prefix conditioning}
% \label{sec:method-arch}

% We fine-tune a single BARTpho-syllable model \citep{tran2022bartpho}---an
% mBART-architecture \citep{liu2020mbart} encoder--decoder operating on syllable
% tokens---to perform all five task variants. Following the text-to-text recipe
% \citep{raffel2020t5}, the task is selected by a textual prefix: the input string
% is ``\texttt{\{task\}: \{source\}}'', and the decoder produces the task-appropriate
% target per Eq.~\ref{eq:seq2seq}. No task-specific parameters, embeddings, or heads
% are added, so cross-task transfer happens through the shared encoder, decoder, and
% tokenizer. We hypothesize positive transfer in particular between the three
% orthographically noisy input tasks (the two dialect directions and lexical
% normalization), which all confront the same non-standard spellings.

% \subsection{Reverse-corpus construction for Feature C}
% \label{sec:method-reverse}

% ViDia2Std is distributed as (dialect, standard) pairs for normalization. To create
% the novel standard$\rightarrow$dialect task we \emph{swap} source and target within
% each pair and tag the new example with a region-specific prefix derived from the
% pair's annotator-assigned region, e.g.\ a Central pair becomes a
% Central standard-to-dialect example whose source is the standard sentence and
% whose target is the original Central dialect sentence. The reverse corpus is thus
% genuinely parallel and human-authored on the target side; only the task direction
% is synthetic. This is what makes region-conditioned generation learnable from an
% existing normalization resource.

% \subsection{Sampling-based back-translation for data-starved regions}
% \label{sec:method-bt}

% Because the reverse corpus inherits ViDia2Std's regional imbalance, the Northern
% and (to a lesser extent) Southern std$\rightarrow$dialect tasks are data-starved.
% We augment them by back-translation \citep{sennrich2016improving,edunov2018understanding}:
% we take \emph{real} dialect sentences as targets and use the strong
% dialect$\rightarrow$standard model (Feature~A) to generate synthetic standard
% \emph{sources}, yielding additional (synthetic-standard, real-dialect) training
% pairs for Feature~C. Two design choices follow \citet{edunov2018understanding}
% directly. First, we \emph{sample} the synthetic sources (temperature with top-$k$
% and top-$p$ truncation) rather than beam-decode them, to avoid the overly regular
% sources that cause mode collapse. Second, we keep the dialect target untouched, so
% the model is never trained to imitate machine output. Synthetic pairs pass a
% light quality filter (round-trip agreement, a length-ratio bound, and a minimum of
% one dialectal marker) before being added. Augmentation is concentrated on Northern
% and Southern, the regions where the marginal value of additional data is highest.

% \subsection{Region rebalancing}
% \label{sec:method-rebalance}

% In addition to back-translation, we physically \emph{oversample} the
% under-represented regions when assembling the training set, so that the three
% std$\rightarrow$dialect tasks are mixed in more even proportions than the raw
% corpus would give. Oversampling is applied at data-build time (duplicating
% records) rather than via a sampler, which keeps the training pipeline simple and
% the effective epoch length explicit.

% \subsection{Transfer-aware model selection}
% \label{sec:method-select}

% A naive selection criterion---validation BLEU pooled over all tasks---would reward
% the copying behaviour described in Section~\ref{sec:bg-one2many}, because copying
% scores well on the high-overlap dialect tasks. We instead select the checkpoint by
% a \emph{transfer-aware} criterion: the macro-average of \emph{per-task} validation
% quality, combining normalization BLEU with the Dialectal Edit-Recall of
% Section~\ref{sec:metrics} for the generation tasks, so that a checkpoint is only
% preferred if it actually performs the dialectal substitutions rather than copying
% the input.

% =============================================================================
\section{Method}
\label{sec:method}

\subsection{Architecture and multi-task task-prefix conditioning}
\label{sec:method-arch}

We fine-tune a single BARTpho-syllable model \citep{tran2022bartpho}---an
mBART-style encoder--decoder \citep{liu2020mbart} operating on Vietnamese
syllable tokens---for all tasks in a unified text-to-text format. Each example is
associated with a task condition \(c\), indicating one of five mappings:
dialect-to-standard normalization, lexical normalization, or
standard-to-dialect generation for the Northern, Central, and Southern regions.
The implementation realizes this condition as a short textual prefix prepended to
the source sentence, following the text-to-text recipe of \citet{raffel2020t5}.

For an input sentence \(x=(x_1,\dots,x_m)\), target sentence
\(y=(y_1,\dots,y_n)\), and task condition \(c\), the encoder receives the
prefixed source
\begin{equation}
  \tilde{x} = [p_c;\,x],
  \label{eq:prefixed-source}
\end{equation}
where \(p_c\) is the prefix associated with task \(c\). The decoder then defines
the conditional distribution
\begin{equation}
  P_\theta(y\mid x,c)
  =
  \prod_{t=1}^{n}
  P_\theta\!\left(y_t \mid y_{<t}, \tilde{x}\right).
  \label{eq:method-seq2seq}
\end{equation}
No task-specific parameters, embeddings, or prediction heads are introduced.
Thus, the same encoder, decoder, and tokenizer are shared across all tasks, while
the prefix provides the only signal that distinguishes the desired mapping.

Given a multi-task training set \(\widetilde{\mathcal{D}}\), the model is trained
with the teacher-forced negative log-likelihood
\begin{equation}
  \mathcal{L}(\theta)
  =
  -\frac{1}{|\widetilde{\mathcal{D}}|}
  \sum_{(x,y,c)\in\widetilde{\mathcal{D}}}
  \sum_{t=1}^{|y|}
  \log P_\theta\!\left(y_t \mid y_{<t}, \tilde{x}\right).
  \label{eq:method-loss}
\end{equation}
This objective is identical across the three features; only the prefix and the
source--target direction change.

\subsection{Reverse-corpus construction for Feature C}
\label{sec:method-reverse}

ViDia2Std is originally distributed as dialect-to-standard parallel data. We
write each example as \((d_i,s_i,r_i)\), where \(d_i\) is the dialect sentence,
\(s_i\) is the corresponding standard sentence, and \(r_i\) is the annotated
dialect region. The normalization direction uses the pair in its original order:
\begin{equation}
  d_i \longrightarrow s_i.
  \label{eq:forward-pair}
\end{equation}
To construct Feature~C, we reverse the same pair and condition the model on the
region label:
\begin{equation}
  s_i \longrightarrow d_i
  \qquad \text{conditioned on } r_i.
  \label{eq:reverse-pair}
\end{equation}
Thus, the standard sentence becomes the source, the original human-written
dialect sentence remains the target, and the region determines which
standard-to-dialect prefix is attached to the input. This gives a genuinely
parallel training signal for region-conditioned generation without introducing
machine-generated dialect targets.

Lexical normalization is added in the same text-to-text format. Its examples
map noisy social-media or teencode input to standard Vietnamese:
\begin{equation}
  x_j^{\mathrm{noise}} \longrightarrow y_j^{\mathrm{std}}.
  \label{eq:lexnorm-pair}
\end{equation}
All three features can therefore be optimized with the same sequence-to-sequence
objective in Eq.~\ref{eq:method-loss}.

\subsection{Sampling-based back-translation for data-starved regions}
\label{sec:method-bt}

The reversed ViDia2Std corpus inherits the original regional imbalance: Central
has many more examples than Southern and especially Northern. To increase the
amount of training data for the data-starved regions, we apply
back-translation \citep{sennrich2016improving,edunov2018understanding}.

For a real dialect sentence \(d\) from region \(r\), we first use the
dialect-to-standard model to generate a synthetic standard source \(\hat{s}\):
\begin{equation}
  \hat{s}
  \sim
  P_{\theta_A}(s \mid d),
  \label{eq:bt-sample}
\end{equation}
where \(\theta_A\) denotes the normalization model. We then form an additional
standard-to-dialect training pair
\begin{equation}
  \hat{s} \longrightarrow d
  \qquad \text{conditioned on } r.
  \label{eq:bt-pair}
\end{equation}
The important asymmetry is that only the source side is synthetic; the target
side remains real dialectal text. This prevents the generator from learning to
imitate machine-produced dialect.

Following \citet{edunov2018understanding}, we sample the synthetic standard
source rather than relying on deterministic beam search. Sampling produces more
varied sources and reduces the risk of training on overly regular synthetic
inputs. A generated pair is retained only if it passes the filtering function
\begin{equation}
  Q(\hat{s},d,r)
  =
  Q_{\mathrm{rt}}(\hat{s},d)
  \cdot
  Q_{\mathrm{len}}(\hat{s},d)
  \cdot
  Q_{\mathrm{mark}}(d,r),
  \label{eq:bt-filter}
\end{equation}
where \(Q_{\mathrm{rt}}\) denotes round-trip agreement,
\(Q_{\mathrm{len}}\) enforces a length-ratio constraint, and
\(Q_{\mathrm{mark}}\) checks that the real dialect target contains at least one
marker associated with the intended region. The filtered back-translated set for
region \(r\) is
\begin{equation}
  \mathcal{B}_r
  =
  \{(\hat{s},d,r): Q(\hat{s},d,r)=1\}.
  \label{eq:bt-set}
\end{equation}
This augmentation is concentrated on Northern and Southern, where the additional
training signal is most useful.

\subsection{Region rebalancing}
\label{sec:method-rebalance}

Back-translation is combined with explicit rebalancing of the standard-to-dialect
training data. Let \(\mathcal{C}_r\) be the original reversed corpus for region
\(r\), and let \(\alpha_r\) be its oversampling factor. The effective
standard-to-dialect training set is
\begin{equation}
  \widetilde{\mathcal{C}}
  =
  \bigcup_{r\in\{\mathrm{N},\mathrm{C},\mathrm{S}\}}
  \left(
    \alpha_r \mathcal{C}_r
    \cup
    \mathcal{B}_r
  \right),
  \label{eq:rebalanced-corpus}
\end{equation}
where \(\alpha_r \mathcal{C}_r\) denotes repeated copies of examples from region
\(r\). Regions with fewer original examples receive larger effective weight. The
full multi-task training set is then
\begin{equation}
  \widetilde{\mathcal{D}}
  =
  \mathcal{D}_{\mathrm{norm}}
  \cup
  \mathcal{D}_{\mathrm{lex}}
  \cup
  \widetilde{\mathcal{C}},
  \label{eq:full-training-set}
\end{equation}
where \(\mathcal{D}_{\mathrm{norm}}\) is the dialect-to-standard normalization
data and \(\mathcal{D}_{\mathrm{lex}}\) is the lexical-normalization data.
Oversampling is applied at data-build time rather than through a dynamic sampler,
which keeps the effective epoch length explicit and the training mixture
reproducible.

\subsection{Transfer-aware model selection}
\label{sec:method-select}

Validation BLEU alone is not a reliable selection criterion for
standard-to-dialect generation. Since many dialectal rewrites differ from the
standard input by only a few tokens, a model that copies the source can obtain
high overlap scores while making little or no dialectal transfer. We therefore
select checkpoints using a transfer-aware validation score.

Let \(S_c(\theta)\) be the validation score for task condition \(c\). For
normalization tasks we use the usual overlap-based score, while for
standard-to-dialect generation we use Dialectal Edit-Recall:
\begin{equation}
  S_c(\theta)
  =
  \begin{cases}
    \mathrm{BLEU}_{c}(\theta)/100,
    & \text{if } c \text{ is a normalization task}, \\[3pt]
    \mathrm{EditRecall}_{c}(\theta),
    & \text{if } c \text{ is a standard-to-dialect task}.
  \end{cases}
  \label{eq:selection-score}
\end{equation}
The selected checkpoint is
\begin{equation}
  \theta^\star
  =
  \operatorname*{arg\,max}_{\theta\in\mathcal{C}}
  \frac{1}{|\mathcal{G}|}
  \sum_{c\in\mathcal{G}} S_c(\theta),
  \label{eq:selected-checkpoint}
\end{equation}
where \(\mathcal{C}\) is the set of saved checkpoints and \(\mathcal{G}\) is the
set of validation task conditions. This macro-average gives each task equal
influence and prevents the selected checkpoint from being dominated by
high-overlap copying behaviour.



% =============================================================================
\section{Baselines}
\label{sec:baselines}

We compare the learned model against four non-neural baselines, all evaluated with
the identical metric code of Section~\ref{sec:metrics} so that numbers are
directly comparable.

\begin{itemize}[leftmargin=1.4em,itemsep=2pt,topsep=2pt]
  \item \textbf{Leave-As-Is (LAI / copy).} The output equals the input. This is the
    crucial floor for one-to-many tasks: because dialectal rewrites touch few
    tokens, LAI attains high BLEU/chrF despite performing zero transfer (its
    copy-rate is $1.0$ and its Dialectal Edit-Recall is $0$ by construction). Any
    model that does not beat LAI on the \emph{transfer} metrics is not doing the
    task.
  \item \textbf{Dictionary.} A token-level substitution map learned from the
    training pairs (most frequent target token for each source token), applied
    position-wise. It transfers aggressively (copy-rate $0$) but ignores context.
  \item \textbf{Rule-based.} Hand-crafted, per-region substitution rules over the
    curated dialectal markers. This is the strongest non-neural baseline on overlap
    metrics and a useful sanity check on the marker lists.
  \item \textbf{Retrieval.} For a query, return the target of its nearest training
    \emph{source} under SBERT sentence embeddings. This tests how far pure
    memorization of the training set can go.
\end{itemize}

Table~\ref{tab:baselines} reports the four baselines per task. The decisive
reading is the \emph{copy trap}: Leave-As-Is posts a high BLEU on the
high-overlap dialect tasks (e.g.\ $74.98$ on Northern, where the dialect is
orthographically closest to standard), yet its Dialectal Edit-Recall is exactly
$0$---it performs no transfer at all. The Rule-based system is the strongest
non-neural method on overlap but still recalls only a small fraction of the gold
substitutions ($0.10$--$0.29$ across regions), while the Dictionary transfers
aggressively (copy-rate $0$) at a catastrophic cost to fluency (BLEU below $6$)
and the Retrieval baseline is uniformly weak. The improved model
(Section~\ref{sec:results}) beats \emph{every} baseline on Edit-Recall and DFR
\emph{while also} clearing the copy BLEU floor.
Figure~\ref{fig:baseline} visualizes the comparison.

\begin{table}[t]
  \centering
  \small
  \caption{Baseline results per task on the test split (BLEU/chrF on a 0--100
  scale; ROUGE-L, copy-rate, DFR, edit-recall in $[0,1]$; ERR may be negative).
  The \textbf{transfer} column is WER/CER for Dialect-to-standard, DFR/edit-recall
  for the std$\rightarrow$dialect tasks, and ERR for Lexical normalization. Note that
  Leave-As-Is (copy) scores high BLEU yet edit-recall $0$---the copy trap.}
  \label{tab:baselines}
  \resizebox{\linewidth}{!}{%
  \begin{tabular}{l l c c c c c}
    \toprule
    \textbf{Task} & \textbf{Baseline} & \textbf{BLEU} & \textbf{chrF} & \textbf{ROUGE-L} & \textbf{copy} & \textbf{transfer}\\
    \midrule
    \multirow{4}{*}{Dialect-to-standard}
      & LAI        & 48.55 & 63.88 & 0.770 & 1.000 & WER 0.321 / CER 0.195\\
      & Dictionary &  3.65 & 23.31 & 0.554 & 0.000 & WER 0.683 / CER 0.532\\
      & Rule-based & 55.29 & 69.03 & 0.828 & 0.356 & WER 0.277 / CER 0.163\\
      & Retrieval  &  3.99 & 14.54 & 0.310 & 0.000 & WER 1.024 / CER 0.798\\
    \midrule
    \multirow{4}{*}{Central standard-to-dialect}
      & LAI        & 44.23 & 61.90 & 0.739 & 1.000 & DFR 0.308 / edit 0.000\\
      & Dictionary &  4.45 & 24.42 & 0.586 & 0.000 & DFR 0.495 / edit 0.432\\
      & Rule-based & 50.40 & 67.74 & 0.810 & 0.233 & DFR 0.651 / edit 0.291\\
      & Retrieval  &  4.00 & 14.63 & 0.312 & 0.000 & DFR 0.305 / edit 0.200\\
    \midrule
    \multirow{4}{*}{Southern standard-to-dialect}
      & LAI        & 47.91 & 65.90 & 0.763 & 1.000 & DFR 0.598 / edit 0.000\\
      & Dictionary &  5.14 & 25.03 & 0.547 & 0.000 & DFR 0.508 / edit 0.435\\
      & Rule-based & 50.54 & 67.38 & 0.776 & 0.555 & DFR 0.677 / edit 0.104\\
      & Retrieval  &  5.93 & 16.15 & 0.316 & 0.000 & DFR 0.365 / edit 0.192\\
    \midrule
    \multirow{4}{*}{Northern standard-to-dialect}
      & LAI        & 74.98 & 86.43 & 0.942 & 1.000 & DFR 0.740 / edit 0.000\\
      & Dictionary & 23.20 & 44.25 & 0.718 & 0.000 & DFR 0.707 / edit 0.468\\
      & Rule-based & 76.73 & 87.26 & 0.944 & 0.537 & DFR 0.890 / edit 0.154\\
      & Retrieval  &  2.66 & 13.26 & 0.286 & 0.000 & DFR 0.455 / edit 0.148\\
    \midrule
    \multirow{4}{*}{Lexical normalization}
      & LAI        & 68.29 & 78.73 & 0.847 & 1.000 & ERR $0.000$\\
      & Dictionary &  6.92 & 26.48 & 0.545 & 0.000 & ERR $-2.088$\\
      & Rule-based & 68.29 & 78.73 & 0.847 & 1.000 & ERR $0.000$\\
      & Retrieval  &  4.93 & 14.24 & 0.290 & 0.000 & ERR $-2.781$\\
    \bottomrule
  \end{tabular}%
  }
\end{table}





\begin{figure}[t!]
  \centering
  \includegraphics[width=0.85\linewidth]{\figures/baseline_comparison.pdf}
  \caption{Baseline comparison across tasks. The contrast between high overlap and
  low transfer for the copy baseline is the key reading.
   }
  \label{fig:baseline}
\end{figure}

% =============================================================================
\section{Evaluation Metrics}
\label{sec:metrics}

We use overlap and semantic metrics where they are valid (normalization) and add
transfer-specific metrics where they are not (generation). Unless noted, BLEU and
chrF are reported on a $0$--$100$ scale and the remaining quantities in $[0,1]$.

\paragraph{BLEU.}
Corpus BLEU \citep{papineni2002bleu} is the geometric mean of modified
$n$-gram precisions ($n=1\!\dots\!4$) times a brevity penalty $BP$:
\begin{equation}
  \mathrm{BLEU} = BP \cdot \exp\!\Big(\textstyle\sum_{n=1}^{4} w_n \log p_n\Big),
  \qquad
  BP = \min\!\big(1,\, e^{1 - r/h}\big),
\end{equation}
with $p_n$ the clipped $n$-gram precision, $w_n=\tfrac14$, $h$ the hypothesis
length and $r$ the reference length. We compute it with SacreBLEU
\citep{post2018sacrebleu} for reproducibility. It is appropriate for the
many-to-one normalization direction but weak for one-to-many generation
(Section~\ref{sec:bg-one2many}).

\paragraph{chrF.}
chrF \citep{popovic2015chrf} is an $F$-score over character $n$-grams (we use the
chrF++ variant that also includes low-order word $n$-grams). Operating on
characters makes it robust to tokenization and to the rich diacritics of
Vietnamese.

\paragraph{ROUGE-L.}
ROUGE-L \citep{lin2004rouge} is the $F$-measure of the longest common subsequence
between hypothesis and reference, rewarding in-order content overlap without
requiring contiguity.

\paragraph{METEOR.}
METEOR \citep{banerjee2005meteor} aligns unigrams (exact and stem matches) and
computes a recall-weighted harmonic mean with a fragmentation penalty,
correlating better with human judgement than raw precision metrics.

\paragraph{WER / CER.}
For comparability with ViDia2Std \citep{ta2026vidia2std} on the normalization
direction we report Word and Character Error Rate---the Levenshtein edit distance
between hypothesis and reference, normalized by reference length (lower is better):
\begin{equation}
  \mathrm{WER} = \frac{S + D + I}{N_{\mathrm{ref}}^{\mathrm{words}}},
  \qquad
  \mathrm{CER} = \frac{S_c + D_c + I_c}{N_{\mathrm{ref}}^{\mathrm{chars}}},
\end{equation}
with $S,D,I$ the substitution, deletion, and insertion counts.

\paragraph{BERTScore.}
BERTScore \citep{zhang2020bertscore} matches hypothesis and reference tokens by
contextual-embedding cosine similarity and reports the $F_1$ of the greedy
matching, capturing semantic equivalence that surface metrics miss; we use a
multilingual model for Vietnamese.

\paragraph{ERR (Error Reduction Rate) --- Feature B.}
For lexical normalization we report the error reduction over the do-nothing (LAI)
baseline. Letting $E(\cdot,\cdot)$ count positionally mismatched tokens against the
reference,
\begin{equation}
  \mathrm{ERR} \;=\; 1 - \frac{E(\hat{y}, y)}{E(x, y)},
  \label{eq:err}
\end{equation}
where $x$ is the noisy source, $\hat{y}$ the system output, and $y$ the gold
normalization. $\mathrm{ERR}=1$ means all normalization errors were fixed,
$\mathrm{ERR}=0$ matches LAI, and $\mathrm{ERR}<0$ means the system introduced more
errors than it fixed.

\paragraph{DFR (Dialect Feature Recall) --- Feature C.}
DFR measures whether the output contains the region's characteristic dialectal
markers that the reference exhibits. Let $M_r$ be the curated marker set for region
$r$, and for an example let $R = \{m \in M_r : m \in y\}$ be the markers present in
the reference. The per-example score is the fraction of those markers reproduced by
the system output $\hat{y}$:
\begin{equation}
  \mathrm{DFR} \;=\; \frac{\bigl|\{m \in R : m \in \hat{y}\}\bigr|}{|R|},
  \qquad |R| > 0,
  \label{eq:dfr}
\end{equation}
averaged over examples whose reference contains at least one marker, and reported
both per region and overall. DFR directly answers ``does it sound like the right
region?'' but depends on a marker list, so we pair it with a marker-independent
measure.

\paragraph{Dialectal Edit-Recall (marker-independent).}
This metric asks whether the model makes the \emph{same token changes} that the
gold reference makes relative to the standard source, without reference to any
marker list. Let $T(\cdot)$ be the set of (lowercased) tokens of a sentence. The
tokens the reference \emph{introduces} over the source are
$A_{\mathrm{gold}} = T(y) \setminus T(x)$, and those the model introduces are
$A_{\mathrm{sys}} = T(\hat{y}) \setminus T(x)$. Then
\begin{align}
  \mathrm{Recall} &= \frac{|A_{\mathrm{gold}} \cap A_{\mathrm{sys}}|}{|A_{\mathrm{gold}}|},
  &
  \mathrm{Precision} &= \frac{|A_{\mathrm{gold}} \cap A_{\mathrm{sys}}|}{|A_{\mathrm{sys}}|},
  \label{eq:editrecall}
\end{align}
with $F_1$ their harmonic mean, averaged over examples with $|A_{\mathrm{gold}}|>0$.
Recall captures how many of the intended dialectal substitutions were made;
precision penalizes spurious changes. Because a copy makes no substitutions,
Leave-As-Is has Edit-Recall $0$ regardless of its BLEU, which is precisely why this
metric, and not BLEU, anchors transfer-aware model selection
(Section~\ref{sec:method-select}).

\paragraph{Region accuracy (marker-based).}
To check controllability, we classify each output to the region whose markers it
most contains and report the fraction matching the \emph{requested} region (over
outputs containing at least one marker), together with the fraction of outputs that
contain no marker of any region (a no-marker / under-generation rate).

\paragraph{Copy-rate.}
The fraction of outputs identical (case-insensitively) to the input. High copy-rate
exposes under-generation and, read alongside Edit-Recall, separates genuine
transfer from copying.

% =============================================================================
\section{Results}
\label{sec:results}

We report the final multi-task model on the test split. Table~\ref{tab:main}
collects the headline numbers; Table~\ref{tab:improvement} isolates the gains of
our data and selection strategy over the baseline model; the baselines of
Table~\ref{tab:baselines} provide the non-neural reference points.

\paragraph{Normalization (Feature A) positioning.}
On dialect$\rightarrow$standard our syllable-level model reaches $82.10$ BLEU,
$88.44$ chrF, and $0.941$ ROUGE-L, with WER $0.1120$ and CER $0.0695$. For
context, the ViDia2Std benchmark \citep{ta2026vidia2std} reports its strongest
system, mBART-large-50 (611M parameters), at $0.8166$ BLEU / $0.9384$ ROUGE-L /
$0.1226$ WER / $0.0754$ CER, and their BARTpho-syllable-base (132M) at $0.7597$
BLEU. Our 396M-parameter model thus \emph{matches and slightly edges} the much
larger SOTA system---its WER ($0.1120$) and CER ($0.0695$) are both below
mBART-large-50's---which we present as a strong positioning result with fewer
parameters rather than a claim of new state of the art. Figure~\ref{fig:main}
visualizes the per-task numbers.

\paragraph{Generation (Feature C).}
On standard$\rightarrow$dialect the model attains an overall DFR of about $0.90$
(per-region: Central $0.901$, Southern $0.897$, Northern $0.897$) and a
marker-independent Dialectal Edit-Recall of $0.714$ (Central), $0.646$
(Southern), and $0.635$ (Northern). Crucially the model's copy-rate is far below
the Leave-As-Is floor of $1.0$---$0.005$ on Central, $0.018$ on Southern, and
$0.090$ on Northern (the hardest region)---confirming that the gains come from
genuine transfer rather than copying. Region accuracy is high on Central
($0.917$) but lower on Southern ($0.424$) and Northern ($0.209$); as
Section~\ref{sec:analysis} explains, this marker-based figure is depressed by
overlapping marker inventories rather than by a real failure to transfer.
Figure~\ref{fig:copy} contrasts copy-driven BLEU with model BLEU and
Figure~\ref{fig:dfr} breaks DFR down by region.

\paragraph{Lexical normalization (Feature B).}
On ViLexNorm the model reaches an Error Reduction Rate of $0.593$ and $87.45$
BLEU, exceeding the $0.5774$ ERR of the ViLexNorm BARTpho-syllable result.

\paragraph{Improvement over the baseline model.}
Table~\ref{tab:improvement} contrasts the baseline model (before targeted
rebalancing and back-translation) with the improved model. The gains concentrate
almost entirely on \emph{Northern}, the data-starved region that rebalancing and
back-translation were designed to help: its Edit-Recall rises from $0.385$ to
$0.635$, its copy-rate \emph{falls} from $0.209$ to $0.090$, its BLEU climbs from
$80.2$ to $87.7$, and its DFR improves from $0.817$ to $0.897$. Central and
Southern are maintained or slightly improved---Southern Edit-Recall rises from
$0.617$ to $0.646$ and DFR from $0.889$ to $0.897$, while Central moves within
noise---and Feature~A normalization is preserved (BLEU $81.68\!\rightarrow\!82.10$,
WER $0.1157\!\rightarrow\!0.1120$). The improvement is exactly where the data
intervention was aimed, with no regression elsewhere.

\begin{table}[t]
  \centering
  \small
  \caption{Main results of the multi-task model per task (test split). BLEU/chrF
  on a 0--100 scale; ROUGE-L, copy-rate, DFR, edit-recall in $[0,1]$; WER/CER lower
  is better. ViDia2Std reference numbers are quoted from \citet{ta2026vidia2std}
  (BLEU/ROUGE-L rescaled to the 0--100 / $[0,1]$ conventions used here).}
  \label{tab:main}
  \resizebox{\linewidth}{!}{%
  \begin{tabular}{l c c c c c}
    \toprule
    \textbf{Task} & \textbf{BLEU} & \textbf{chrF} & \textbf{ROUGE-L} & \textbf{copy} & \textbf{task-metric}\\
    \midrule
    Dialect-to-standard (A)            & 82.10 & 88.44 & 0.941 & 0.010 & WER 0.112 / CER 0.069\\
    Central standard-to-dialect (C)   & 69.09 & 80.14 & 0.906 & 0.005 & DFR 0.901 / edit 0.714\\
    Southern standard-to-dialect (C)  & 70.36 & 81.74 & 0.892 & 0.018 & DFR 0.897 / edit 0.646\\
    Northern standard-to-dialect (C)  & 87.68 & 93.47 & 0.977 & 0.090 & DFR 0.897 / edit 0.635\\
    Lexical normalization (B)                & 87.45 & 93.09 & 0.963 & 0.039 & ERR 0.593\\
    \midrule
    \emph{ref:} mBART-large-50 (A) \citep{ta2026vidia2std} & 81.66 & --- & 0.938 & --- & WER 0.123 / CER 0.075\\
    \emph{ref:} BARTpho-syll-base (A) \citep{ta2026vidia2std} & 75.97 & --- & --- & --- & ---\\
    \bottomrule
  \end{tabular}%
  }
\end{table}

\begin{table}[t]
  \centering
  \small
  \caption{Improvement of the final model over the baseline model on the four
  dialect tasks (test split), \emph{baseline$\rightarrow$improved}. For
  Dialect-to-standard the transfer column is WER/CER (lower is better); for the
  std$\rightarrow$dialect tasks it is DFR / Edit-Recall. The headline is
  \textbf{Northern}, the data-starved region targeted by rebalancing and
  back-translation: Edit-Recall and DFR rise sharply while copy-rate \emph{falls}.
  Central and Southern are maintained or slightly improved, and Feature~A is
  preserved.}
  \label{tab:improvement}
  \begin{tabular}{l c c c c}
    \toprule
    \textbf{Task} & \textbf{BLEU} & \textbf{DFR / WER} & \textbf{Edit-Rec / CER} & \textbf{copy-rate}\\
    \midrule
    Dialect-to-standard (A)
      & $81.68\!\rightarrow\!82.10$ & $0.1157\!\rightarrow\!0.1120$ & $0.0713\!\rightarrow\!0.0695$ & $0.010\!\rightarrow\!0.010$\\
    Central standard-to-dialect (C)
      & $69.51\!\rightarrow\!69.09$ & $0.896\!\rightarrow\!0.901$ & $0.719\!\rightarrow\!0.714$ & $0.008\!\rightarrow\!0.005$\\
    Southern standard-to-dialect (C)
      & $69.06\!\rightarrow\!70.36$ & $0.889\!\rightarrow\!0.897$ & $0.617\!\rightarrow\!0.646$ & $0.024\!\rightarrow\!0.018$\\
    \textbf{Northern standard-to-dialect} (C)
      & $\mathbf{80.16\!\rightarrow\!87.68}$ & $\mathbf{0.817\!\rightarrow\!0.897}$ & $\mathbf{0.385\!\rightarrow\!0.635}$ & $\mathbf{0.209\!\rightarrow\!0.090}$\\
    \bottomrule
  \end{tabular}
\end{table}

\begin{figure}[t]
  \centering
  \includegraphics[width=0.85\linewidth]{\figures/main_results.pdf}
  \caption{Main multi-task results across the five task variants.
   }
  \label{fig:main}
\end{figure}

\begin{figure}[t]
  \centering
  \begin{subfigure}{0.48\linewidth}
    \centering
    \includegraphics[width=\linewidth]{\figures/copy_vs_model_bleu.pdf}
    \caption{Copy (LAI) BLEU vs.\ model BLEU.}
    \label{fig:copy}
  \end{subfigure}\hfill
  \begin{subfigure}{0.48\linewidth}
    \centering
    \includegraphics[width=\linewidth]{\figures/dfr_by_region.pdf}
    \caption{Dialect Feature Recall by region.}
    \label{fig:dfr}
  \end{subfigure}
  \caption{The copy baseline inflates BLEU (left), motivating the transfer metrics
  (right). \emph{(Figures generated by a separate plotting step.)}}
  \label{fig:copy-dfr}
\end{figure}

% =============================================================================
\section{Analysis and Error Taxonomy}
\label{sec:analysis}

We organize the analysis around five error modes of the generation direction and
the metric that detects each: (i) \emph{under-transfer / copying} (copy-rate,
Edit-Recall); (ii) \emph{wrong-region leakage} (region accuracy / the confusion
matrix); (iii) \emph{over-transfer} (Edit-Precision); (iv) \emph{semantic drift}
(BERTScore); and (v) \emph{regional sparsity effects}, concentrated on Northern.
The remainder of this section grounds each mode in the model's actual outputs.

\subsection{Learned lexical substitutions}
\label{sec:analysis-subs}

Table~\ref{tab:subs} lists the most frequent standard$\rightarrow$dialect token
substitutions the model produced on the test set, per region (mined directly from
output/source alignment). These are not hand-written rules: the model
\emph{induced} them from the parallel data. They are linguistically correct,
well-attested regional lexicalisations---Central \emph{này}$\rightarrow$\emph{ni},
\emph{vậy}$\rightarrow$\emph{rứa}, \emph{đâu}$\rightarrow$\emph{mô},
\emph{gì}$\rightarrow$\emph{chi}, \emph{sao}$\rightarrow$\emph{răng},
\emph{tôi}$\rightarrow$\emph{tui}; Northern orthographic/lexical variants
\emph{trời}$\rightarrow$\emph{giời}, \emph{mẹ}$\rightarrow$\emph{u}/\emph{bu},
\emph{nhỉ}$\rightarrow$\emph{nhở}, \emph{chưa}$\rightarrow$\emph{chửa}; and Southern
\emph{về}$\rightarrow$\emph{dìa}, \emph{không}$\rightarrow$\emph{hông},
\emph{thật}$\rightarrow$\emph{thiệt}, \emph{xem}$\rightarrow$\emph{coi},
\emph{đắt}$\rightarrow$\emph{mắc}. That the model recovers the canonical function
words, pronouns, and kinship terms of each region from data alone is direct
evidence that Feature~C learns genuine dialect lexicon rather than surface noise.

\begin{table}[t]
  \centering
  \small
  \caption{Most frequent learned standard$\rightarrow$dialect lexical
  substitutions per region (token, with output count), mined from the model's test
  outputs. All are well-attested regional forms.}
  \label{tab:subs}
  \begin{tabular}{l l l l}
    \toprule
    & \textbf{Central} & \textbf{Northern} & \textbf{Southern}\\
    \midrule
    1 & này $\rightarrow$ ni (15)   & nhỉ $\rightarrow$ nhờ (28)   & về $\rightarrow$ dìa (5)\\
    2 & vậy $\rightarrow$ rứa (10)  & nhỉ $\rightarrow$ nhể (18)   & không $\rightarrow$ hông (5)\\
    3 & nào $\rightarrow$ mô (5)    & nhỉ $\rightarrow$ nhở (12)   & thật $\rightarrow$ thiệt (3)\\
    4 & gì $\rightarrow$ chi (4)    & giàu $\rightarrow$ giầu (11) & xem $\rightarrow$ coi (3)\\
    5 & chứ $\rightarrow$ chơ (4)   & trời $\rightarrow$ giời (10) & đắt $\rightarrow$ mắc (3)\\
    6 & đâu $\rightarrow$ mô (4)    & chưa $\rightarrow$ chửa (9)  & này $\rightarrow$ nè (1)\\
    7 & sao $\rightarrow$ răng (3)  & mẹ $\rightarrow$ u (5)       & chăn $\rightarrow$ mền (1)\\
    8 & tôi $\rightarrow$ tui (2)   & mẹ $\rightarrow$ bu (4)      & hạt $\rightarrow$ hột (1)\\
    \bottomrule
  \end{tabular}
\end{table}

\subsection{Region controllability and the confusion matrix}
\label{sec:analysis-region}

Table~\ref{tab:confusion} gives the region-confusion matrix: rows are the
\emph{requested} region, columns are the region whose markers a heuristic detector
found in the output (plus a ``none'' column for outputs with no detected marker).
Central is well controlled---$850$ of $886$ marked outputs land on the diagonal.
Northern and Southern leak: Northern requests are frequently detected as Central
($79$) and Southern requests split across Central ($96$) and Northern ($98$). This
is partly genuine confusion under imbalance, but it is substantially an
\emph{artefact of overlapping marker inventories}: several markers (e.g.\
\emph{mô}, \emph{tui}) belong to more than one region, and Northern dialect is
orthographically the closest to standard, so it carries few distinctive markers
for the detector to fire on. \textbf{For honesty we therefore do not report a raw
``wrong-region'' error count}---the heuristic marker classifier inflates it through
marker overlap---and instead report this confusion matrix together with the
marker-independent Dialectal Edit-Recall, which credits the model for making the
correct token changes regardless of which region's marker list claims them.

\begin{table}[t]
  \centering
  \small
  \caption{Region-confusion matrix. Rows: requested region; columns: region
  detected by the heuristic marker classifier in the output (``none'' = no marker
  found). Off-diagonal mass is inflated by overlapping marker inventories and by
  Northern's orthographic closeness to standard; see text.}
  \label{tab:confusion}
  \begin{tabular}{l c c c c}
    \toprule
    \textbf{Requested $\downarrow$ / Detected $\rightarrow$} & \textbf{Northern} & \textbf{Central} & \textbf{Southern} & \textbf{none}\\
    \midrule
    Northern & 91  & 79  & 16  & 15\\
    Central  & 167 & 850 & 39  & 10\\
    Southern & 98  & 96  & 126 & 10\\
    \bottomrule
  \end{tabular}
\end{table}

\subsection{One-to-many diversity under sampling}
\label{sec:analysis-diversity}

To confirm that Feature~C produces genuine variation---not a single memorized
rewrite---we sample five outputs per standard input and measure Self-BLEU among
them (lower means more diverse). Self-BLEU under sampling is $81.6$ for Central,
$86.7$ for Southern, and $93.7$ for Northern. Sampling thus yields real
variation, most for Central and least for Northern, consistent with Northern
being orthographically closest to standard (fewer alternative realizations to
sample) and with its data scarcity. For instance, the standard ``Tại sao bạn
không trả lời tin nhắn của tôi?'' yields Central samples ranging over
``Mắc mớ gì\dots'', ``Răng bạn\dots'', and ``Mắc chi\dots''---distinct, idiomatic
realizations that a single-reference BLEU would penalize.

\subsection{Tokenization: syllables vs.\ subwords}
\label{sec:analysis-tok}

The probe in Table~\ref{tab:tok} makes Section~\ref{sec:bg-tok}'s robustness
argument concrete. Standard syllables tokenize as whole units
(\texttt{\spm Anh}, \texttt{\spm đi}, \texttt{\spm đâu}, \texttt{\spm rồi}), whereas rarer
dialectal forms decompose into subwords:
\emph{Bả}$\rightarrow$\texttt{[\spm B, ả]},
\emph{rứa}$\rightarrow$\texttt{[\spm r, ứ, a]}, and
\emph{Hông}$\rightarrow$\texttt{[\spm H, ông]}. Because every white-space unit remains
a valid token sequence, out-of-vocabulary dialect spellings degrade gracefully
into shorter meaningful pieces rather than triggering a hard segmentation
failure---exactly the property that motivates BARTpho-syllable for this task.

\begin{table}[t]
  \centering
  \small
  \caption{SentencePiece tokenization of standard vs.\ dialect surface forms.
  Standard syllables are whole pieces; rarer dialect forms fragment into subwords
  (\texttt{\spm} marks a word-initial piece).}
  \label{tab:tok}
  \begin{tabular}{l l}
    \toprule
    \textbf{Surface form} & \textbf{SentencePiece pieces}\\
    \midrule
    ``Anh ấy đi đâu rồi?'' (std) & \texttt{\spm Anh \spm ấy \spm đi \spm đâu \spm rồi ?}\\
    ``Ảnh đi mô rồi?'' (dialect) & \texttt{\spm Ảnh \spm đi \spm mô \spm rồi ?}\\
    ``Bà ấy nói gì vậy?'' (std)  & \texttt{\spm Bà \spm ấy \spm nói \spm gì \spm vậy ?}\\
    ``Bả nói chi rứa?'' (dialect)& \texttt{\spm B ả \spm nói \spm chi \spm r ứ a ?}\\
    ``Không biết'' (std)         & \texttt{\spm Không \spm biết}\\
    ``Hông biết'' (dialect)      & \texttt{\spm H ông \spm biết}\\
    \bottomrule
  \end{tabular}
\end{table}

\subsection{Qualitative examples: Northern before and after}
\label{sec:analysis-qual}

Table~\ref{tab:qual} shows Northern outputs of the baseline model versus the
improved model on the same standard inputs. Where the baseline frequently copied
the standard sentence verbatim, the improved model now produces Northern-specific
forms: ``Tại sao bạn không trả lời\dots'' becomes ``\dots không \emph{giả nhời}\dots''
(the Northern realization of \emph{trả lời}), and ``Mẹ tôi\dots'' becomes
``\emph{u} tôi\dots'' (the Northern kinship term). We note honestly that Northern
remains the hardest region: several inputs (e.g.\ ``Hôm nay trời đẹp quá\dots'',
``Anh ấy nói cái gì vậy\dots'') are still copied unchanged, consistent with its
$0.090$ residual copy-rate and its data scarcity.

\begin{table}[t]
  \centering
  \small
  \caption{Northern std$\rightarrow$dialect: baseline (old) vs.\ improved (new)
  outputs. The improved model introduces Northern forms (\emph{giả nhời},
  \emph{u}) where the baseline copied; the last row shows a still-copied input.}
  \label{tab:qual}
  \begin{tabular}{p{0.30\linewidth} p{0.30\linewidth} p{0.30\linewidth}}
    \toprule
    \textbf{Standard input} & \textbf{Baseline (old)} & \textbf{Improved (new)}\\
    \midrule
    Tại sao bạn không trả lời tin nhắn của tôi?
      & Mắc gì bạn không trả lời tin nhắn của tôi?
      & Tại sao bạn không \emph{giả nhời} tin nhắn của tôi?\\
    Mẹ tôi nấu ăn rất ngon.
      & Bu tôi nấu ăn rất ngon.
      & \emph{u} tôi nấu ăn rất ngon.\\
    Hôm nay trời đẹp quá, chúng ta đi chơi đi.
      & Hôm nay trời đẹp quá, chúng ta đi chơi đi. \emph{(copied)}
      & Hôm nay trời đẹp quá, chúng ta đi chơi đi. \emph{(still copied)}\\
    \bottomrule
  \end{tabular}
\end{table}

A complementary human evaluation of adequacy, fluency, and dialect authenticity
is left as a planned human evaluation to validate the automatic transfer metrics.

% =============================================================================
\section{Limitations}
\label{sec:limitations}

Several limitations bound our claims. \emph{Single reference}: ViDia2Std supplies
one dialectal realization per sentence, so even our transfer-aware metrics are
computed against a single target; the human-eval protocol partly mitigates this but
does not eliminate it. \emph{Marker dependence}: DFR and region accuracy rely on a
curated marker list and therefore inherit its coverage gaps, which is exactly why
we report the marker-independent Edit-Recall alongside them. \emph{Region
granularity}: collapsing 63 provinces into three regions hides within-region
variation. \emph{Data imbalance}: despite oversampling and targeted
back-translation, Northern remains the smallest region and its results are the
least reliable. \emph{Synthetic sources}: back-translated standard sources are
model output and may carry artefacts, although sampling and keeping real dialect
targets limit the damage. \emph{Annotation provenance}: the corpus is drawn from
public Facebook comments, with the attendant register, topic, and demographic
biases.

% =============================================================================
\section{Conclusion}
\label{sec:conclusion}

We presented a single multi-task BARTpho-syllable model that performs
dialect$\rightarrow$standard normalization, lexical normalization, and---as the
novel contribution---region-conditioned standard$\rightarrow$dialect generation
for the three Vietnamese dialect regions. The work delivers three concrete
outcomes. \textbf{(A)} On the well-studied normalization direction our 396M-parameter
model is competitive with a 611M-parameter SOTA system, reaching $82.10$ BLEU at a
WER/CER of $0.1120$/$0.0695$---both below ViDia2Std's mBART-large-50---so we match
the larger model with roughly a third fewer parameters. \textbf{(B)} We built, to
our knowledge, the first region-conditioned standard$\rightarrow$dialect generator
for Vietnamese, reaching an overall DFR of about $0.90$ and Edit-Recall of
$0.635$--$0.714$; targeted rebalancing and sampling back-translation recover most
of the Northern gap (Edit-Recall $0.385\!\rightarrow\!0.635$, copy-rate
$0.209\!\rightarrow\!0.090$) without regressing the other regions. \textbf{(C)} We
introduced transfer-aware metrics---Dialect Feature Recall and a marker-independent
Dialectal Edit-Recall, paired with copy-rate---that expose the copying behaviour a
single-reference BLEU hides and that anchor model selection toward genuine
transfer. Together these establish standard$\rightarrow$dialect generation as a
learnable, measurable task for Vietnamese.

% =============================================================================
\bibliographystyle{plainnat}
\bibliography{references}

\end{document}
